\documentclass{article}
\usepackage{amsmath, amsthm, amsfonts, amssymb}
\theoremstyle{plain}
\newtheorem{thm}{Theorem}
\numberwithin{thm}{section}

\theoremstyle{plain}
\newtheorem{prop}{Proposition}
\numberwithin{prop}{section}

\theoremstyle{definition}
\newtheorem{defn}{Definition}
\numberwithin{defn}{section}

\theoremstyle{remark}
\newtheorem*{rem}{Remark}

\newtheorem*{cor}{Corollary}

\numberwithin{equation}{section}

\newcommand{\op}{\prime}
\begin{document}
\title{Elementary Topology of $\mathbb{R}^n$}
\author{Amey Joshi}
\maketitle
\section{Metric spaces}\label{s1}
\begin{defn}[Metric space]\label{s1d1}
Let $X$ be a non-empty set. A mapping $d: X \times X \mapsto \mathbb{R}$ is
a metric if, for all $x, y, z \in X$:
\begin{enumerate}
\item $d(x, y) \ge 0$ and $d(x, y) = 0$ iff $x = y$.
\item $d(x, y) = d(y, x)$.
\item $d(x, y) + d(y, z) \ge d(x, z)$.
\end{enumerate}
The set $X$ together with the metric $d$ is called a metric space and is
denoted by $(X, d)$.
\end{defn}

\begin{defn}[$k$-cell]\label{s1d2}
If $a_i < b_i$ for $i = 1, \ldots, k$ then a $k$-cell is the set 
\[
\{x \in \mathbb{R}^k\;:\; x = (x_1, \ldots, x_k)\;\text{ and } a_i \le x_i \le
b_i\}
\]
\end{defn}

\begin{defn}[Open ball]\label{s1d3}
Let $(X, d)$ be a metric space. An open ball of radius $r > 0$ centred at 
point $x \in X$ is the set $B_r(x) = \{y \in Y \;:\; d(y, x) < r\}$.
\end{defn}

\begin{defn}[Convex set]\label{s1d4}
A set $E \subseteq \mathbb{R}^k$ is convex if $x, y \in E \Rightarrow \lambda x
+ (1 - \lambda)y \in E$ for all $\lambda \in (0, 1)$.
\end{defn}
\begin{rem}
The operation $\lambda x + (1 - \lambda)y$ is defined only on a vector space.
Therefore, the idea of convexity is applicable only for metric spaces with a
vector space structure.
\end{rem}

\begin{prop}\label{s1p1}
$k$-cell is convex.
\end{prop}
\begin{proof}
Let $a_i < b_i, i = 1, \ldots, k$ be such that E = $\{x \in \mathbb{R}^k \;:\;
x = (x_1, \ldots, x_k), a_i \le x_i \le b_i, i = 1, \ldots, k\}$ is a $k$-cell.
Let $x, y \in E$ and $\lambda \in (0, 1)$. If $z = \lambda x + (1 - \lambda y$
so that $z_i = \lambda x_i + (1 - \lambda)y_i$.

Consider
\begin{eqnarray*}
\lambda a_i &\le \lambda x_i \le& \lambda b_i \\
(1 - \lambda) a_i &\le (1 - \lambda)y_i \le & (1 - \lambda)b_i.
\end{eqnarray*}
Adding them, we get $a_i \le z_i \le b_i$ or that $z \in E$.
\end{proof}

\begin{prop}\label{s1p2}
An open ball is convex.
\end{prop}
\begin{proof}
Let $x, y \in B_r(x_0)$. Then $|x - x_0| < r, |y - y_0| < r$ so that if $z =
\lambda x + (1 - \lambda)y$ then $|z - x_0| = |\lambda x + (1 - \lambda)y -
\lambda x_0 - (1 - \lambda)x_0| \le \lambda|x - x_0| + (1 - \lambda)|y - x_0|
< \lambda r + (1 - \lambda)r = r$.
\end{proof}

The following definitions are for a metric space $(X, d)$ and $E \subseteq X$.
\begin{defn}\label{s1d5}
A neighbourhood of $p$ is the set of points in the ball $B_r(p)$ for some $r>0$.
\end{defn}

\begin{defn}\label{s1d6}
A point $p$ is a limit point of $E$ if every neighbourhood of $p$ has a point
$q \ne p$, such that $q \in E$.
\end{defn}

\begin{defn}\label{s1d7}
If $p$ is not a limit point of $E$ then it called an isolated point of $E$.
\end{defn}

\begin{defn}\label{s1d8}
$E$ is closed if it contains all its limit points.
\end{defn}

\begin{defn}\label{s1d9}
A point $p$ is an interior point of $E$ if there exists $r > 0$ such that 
$B_r(p) \subseteq E$.
\end{defn}

\begin{defn}\label{s1d10}
$E$ is open if every point of $E$ is an interior point.
\end{defn}

\begin{defn}\label{s1d11}
The complement of $E$ is the set $E^c = \{p \in X \;;\; p \notin E\}$.
\end{defn}

\begin{defn}\label{s1d12}
$E$ is perfect if every point of it is a its limit point.
\end{defn}

\begin{defn}\label{s1d13}
$E$ is bounded if there exists $M > 0$ such that for all $p, q \in E$, 
$d(p, q) < M$.
\end{defn}

\begin{defn}\label{s1d14}
$E$ is dense in $X$ is every point of $X$ is either a limit point of $E$ or a
point in $E$.
\end{defn}

\begin{prop}\label{s1p3}
A neighbourhood is an open set.
\end{prop}
\begin{proof}
Consider the neighbourhood $B_r(p)$ of a point $p$ and let $q \in B_r(p)$. Then
$d(p, q) < r$. Let $r^\ast = r - d(p, q)$. Consider a point $s \in 
B_{r^\ast}(q)$. For such a point,
\[
d(p, s) \le d(p, q) + d(q, s) < d(p, q) + r^\ast = r.
\]
\end{proof}

\begin{prop}\label{s1p4}
If $p$ is a limit point of a set $E$ then every neighbourhood of $p$ contains
infinitely many points of $E$.
\end{prop}
\begin{proof}
Let $N$ be a neighbourhood of a limit point $p$ of the set $E$. Since $N$ is a
neighbourhood, so is $B_r(p) \cap N$ for all $r > 0$. Since $p$ is a limit
point, $(B_r(p) \cap N) \cap (E - \{p\}) \ne \emptyset$.

Let, if possible, $N$ be a finite set so that $N \cap (E - \{p\}) = \{q_1, 
\ldots, q_m\}$. Define,
\[
\delta = \min_{1 \le i \le m} d(p, q_i).
\]
As $N$ is a finite set, $\delta > 0$. If we choose $r = \delta/2$, the set 
$(B_{\delta/2}(p) \cap N) \cap (E - \{p\})$ is empty, contradicting our 
hypothesis that $p$ is a limit point.
\end{proof}

An immediate consequence of this proposition is:
\begin{prop}\label{s1p5}
A finite set has no limit points.
\end{prop}
\begin{proof}
Let $E$ be a finite set. Let, if possible, $p \in E$ be its limit point. Then
evert neighbourhood of $p$ has infinitely many points of $E$. Since $E$ is
finite, no such point exists.
\end{proof}

Recall that the proposition $A \Rightarrow B$ is true if $A$ is false or $B$
is true. As a result,
\begin{prop}\label{s1p6}
A non-empty finite set is closed.
\end{prop}

\begin{prop}\label{s1p7}
Any metric space $(X, d)$ viewed as a subset of itself is both open and closed.
\end{prop}
\begin{proof}
For any $r > 0, p \in X$, $B_r(p) = \{x \in X \;:\; d(x, p) < r\}$. By 
definition $B_r(p) \subset X$. Since $p$ is arbitary, this implies that $X$
is open.

A set is closed if it contains all its limit points. The limit points are
members of the ambient space, in this case, $X$ itself. Thus $X$ contains all
its limit points and therefore it is a closed set.
\end{proof}

\begin{prop}\label{s1p8}
Let $E_\alpha$ be an artibrary collection of sets. Then
\[
\left(\cup_\alpha E_\alpha\right)^c = \cap_\alpha E_\alpha^c.
\]
\end{prop}
\begin{proof}
Let
\[
x \in \left(\cup_\alpha E_\alpha\right)^c.
\]
Then $x \notin E_\alpha, \;\forall\; \alpha$, that is $x \in E_\alpha^c, 
\;\forall\alpha$ or that $x \in \cap_\alpha E_\alpha^c$. Therefore,
\[
\left(\cup_\alpha E_\alpha\right)^c \subseteq \cap_\alpha E_\alpha^c.
\]

Not suppose, $y \in \cap_\alpha E_\alpha^c$ then $y \in E_\alpha^c$ or $y
\notin E_\alpha$ for all $\alpha$. Therefore, $y \notin \cup_\alpha E_\alpha$
so that 
\[
y \in \left(\cup_\alpha E_\alpha\right)^c
\]
and hence
\[
\cap_\alpha E_\alpha^c \subseteq \left(\cup_\alpha E_\alpha\right)^c.
\]
\end{proof}

\begin{prop}\label{s1e9}
Let $E$ be open then $E^c$ is closed.
\end{prop}
\begin{proof}
Let $x$ be a limit point of $E^c$. Therefore, every neighbourhood of $x$ has
points in $E^c$. That is, no neighbourhood of $x$ will be entirely in $E$.
That is, $x$ is not an interior point of $E$. Since $E$ is open, all its points
are interior points. Therefore, $x \notin E$, or that $x \in E^c$.

Since $x$ was an arbitrary limit point of $E^c$, this means that every limit
point of $E^c$ lies in $E^c$.
\end{proof}

\begin{prop}\label{s1e10}
Let $E^c$ be closed then $E$ is open.
\end{prop}
\begin{proof}
Let $x$ be a limit point of $E$. If $x$ is an interior point then $x \in E$.
otherwise, every neighbourhood of $x$ has points in $E^c$, making $x$ a limit 
point of $E^c$ as well. But $E^c$ is closed. Therefore, $x \in E^c$. Thus, 
every limit point of $E$ is either an interior point of $E$ or limit point of
$E^c$, making $E$ an open set.
\end{proof}

We can combine the above two propositions into
\begin{thm}\label{s1t1}
A set is open iff its complement is closed.
\end{thm}

Since the complement of the complement of a set is the set itself, this theeorem
is equivalent to
\begin{thm}\label{s1t2}
A set is closed iff its complement is open.
\end{thm}

The next few theorems are important in justifying the definition of a 
topological space.
\begin{thm}\label{s1t3}
An arbitrary union of open sets is open.
\end{thm}
\begin{proof}
Let $\{E_\alpha\}$ be a collection of open sets. Therefore every point in the
collection is an interior point. Therefore, every point in $\cup_\alpha 
E_\alpha$ is an interior point.
\end{proof}

\begin{thm}\label{s1t4}
An arbitrary intersection of closed sets is closed.
\end{thm}
\begin{proof}
Let $\{E_\alpha\}$ be a collection of closed sets. Let $p$ be a limit point of
$F = \cap_\alpha E_\alpha$. Then $B_r(p) \cap (F - \{p\}) \ne \varnothing \;
\forall r > 0$. Therefore, $B_r(p) \cap (E_\alpha - \{p\}) \ne \varnothing \;
\forall r > 0, \;\forall \alpha$, or that $p$ is a limit point of every 
$E_\alpha$. Since $E_\alpha$ are closed, $p \in E_\alpha\;\forall \alpha$ or
that $p \in F$.

Since $p$ was an arbitrary limit point of $F$, every limit point of $F$ belongs
to $F$, making it a closed set.
\end{proof}

An alternate proof is:
\begin{proof}
Let $\{E_\alpha\}$ be a collection of closed sets. Then $\{E_\alpha^c\}$ are
all open so that $\cup_\alpha E_\alpha^c$ is open or
\[
\left(\cap_\alpha E_\alpha\right)^c
\]
is open, or $\cap_\alpha E_\alpha$ is closed.
\end{proof}

\begin{thm}\label{s1t5}
An intersection of a finite collection of open sets is open.
\end{thm}
\begin{proof}
Let $\{E_1, \ldots, E_n\}$ be a finite collection of open sets. Let $x \in 
E_1 \cap \ldots \cap E_n$. Then $x \in E_1, \ldots, x \in E_n$. As all of them
are open sets, there are positive numbers $r_1, \ldots, r_n$ such that 
$B_{r_i} \subset E_i$ for $i = 1, \ldots, n$. Let $r = \min(\{r_1, \ldots, 
r_n\})$. The minimum of a finite set of positive numbers is positive. Therefore,
$B_r(x) \subset E_i$ for all $i = 1, \ldots, n$ or that $B_r(x) \subset E_1
\cap \ldots \cap E_n$. 

Since $x$ was an arbitrary point in $E_1 \cap \ldots \cap E_n$, we showed that
every point in $E_1 \cap \ldots \cap E_n$ is an interior point, making it an
open set.
\end{proof}

\begin{thm}\label{s1t6}
A union of a finite collection of closed sets is closed.
\end{thm}
\begin{proof}
Let $\{E_1, \ldots, E_n\}$ be a finite collection of closed sets. Therefore,
$\{E_1^c, \ldots, E_n^c\}$ is a finite collection of open sets. By the previous
therefore, $E_1^c \cap \ldots E_n^c = (E_1 \cup \ldots \cup E_n)^c$ is open or
$E_1 \cup \ldots E_n$ is closed.
\end{proof}

\begin{defn}\label{s1d15}
Let $(X, d)$ be a metric space and $E \subset X$. Let $E^\op$ be the set of all
limit points of $E$. Then the closure of $E$ is defined as $\bar{E} = E \cup
E^\op$.
\end{defn}

\begin{thm}\label{s1t7}
Let $(X, d)$ be a metric space and $E \subset X$. Then,
\begin{enumerate}
\item $\bar{E}$ is closed.
\item $E = \bar{E}$ iff $E$ is closed.
\item $\bar{E} \subset F$ for every closed set $F \subset X$ such that $E 
\subset F$. 
\end{enumerate}
From points (1) and (3), we see that $\bar{E}$ is the smallest closed subset of
$X$ that contains $E$.
\end{thm}
\begin{proof}
By definition, $\bar{E}$ has all limit points of $E$. Therefore, $\bar{E}$ is 
closed. If $E$ is closed then $E^\op \subset E$ and hence $\bar{E} = E$. 
Conversely, if $\bar{E} = E$ then $E^\op \subset E$, that is, all limit points
of $E$ belong to $E$, making it a closed set.

Let $F \subset X$ be a closed set such that $E \subset F$. Let $p$ be a limit
point of $E$ so that $p \in \bar{E}$. Since $E \subset F$, and $F$ is closed, 
$p \in F$. Therefore, $E^\op \subseteq F$. If $q \in E$ is not a limit point of
$E$ then as as $E \subset F$, $q \in F$. Thus any point in $E \cup E^\op$ is a
member of $F$ so that $\bar{E} \subseteq F$.
\end{proof}

A slightly more detailed proof of item (3) is
\begin{proof}
Let $F \subset X$ be a closed set such that $E \subset F$. Let $p$ be a limit
point of $E$. Then every neighbourhood of $p$ has points in $E$ other than $p$.
Since $E \subset F$, every neighbourhood of $p$ has points in $F$ other than
$p$, making $p$ a limit point of $F$. Since $F$ is closed, $p \in F$. Thus,
$E^\op \subset F$. If $q \in E$ is not a limit point of $E$ then since $E 
\subset F$, $q \in F$. Thus ant point in $E \cup E^\op$ is a member of $F$ so
that $\bar{E} \subseteq F$.
\end{proof}

\begin{thm}\label{s1t8}
Let $E$ be a non-empty set of real numbers which is bounded above. Let $y = 
\sup{E}$. Then $y \in \bar{E}$. If $y \in E$ then $E$ is closed.
\end{thm}
\begin{proof}
If $y = \sup E$ then for all $x \in E$, $x \le y$ and if $y^\op \ge x$ then
$y^\op \ge y$. Consider the set $B_\epsilon(y) = \{x \in \mathbb{R}\;:\; |x - y|
< \epsilon\} = \{x \in \mathbb{R} \;:\; x - \epsilon < y < x + \epsilon\}$. As
$y = \sup E$, numbers in $[x - \epsilon, y) \subseteq E$. This is true for all
$\epsilon$. Therefore, $y$ is a limit point of $E$ so that $y \in \bar{E}$. If
$y \in E$ then $E$ is closed.
\end{proof}

\begin{defn}\label{s1d16}
Let $(X, d)$ be a metric space and $E \subset Y \subset X$. Then $E$ is open 
relative to $Y$ if for each $p \in E$, there exists $r_p > 0$ such that if
$d(p, q) < r_p$ and $q \in Y$ then $q \in E$.
\end{defn}
\begin{rem}
If $E$ is open relative to $Y \subset X$ is for any $p \in E$, we can find 
$\epsilon_p > 0$ such that $B_{\epsilon_p} \cap Y \subset E$.
\end{rem}

\begin{thm}\label{s1t9}
Let $(X, d)$ be a metric space and $Y \subset X$. $E \subset Y$ is open relative
to $Y$ iff $E = Y \cap G$ for some open subset $G$ of $X$.
\end{thm}
\begin{proof}
Let $p \in E$. Since $p$ is also a member of $G$, there exists $\epsilon_p > 0$
such that $B_{\epsilon_p}(p) \subset G$. Choose $\epsilon_p$ so that 
$B_{\epsilon_p} \subset Y$. Then, for a $q \in Y$ such that $q \in 
B_{\epsilon_p}$, $q$ belongs to both $Y$ and $G$ so that it belongs to $E = Y
\cap G$, making $E$ open relative to $Y$.

Now suppose $E$ is open relative to $Y$. Therefore, for a $p \in E$, if we can
find $\epsilon_p > 0$ such that if $d(p, q) < \epsilon_p$ for a point $q \in Y$,
then $q \in E$. We have enclosed a point $p$ of $E$ in an open ball of radius
$\epsilon_p$ such that all $q \in Y$ lying in the ball, also belong to $E$. 
Define
\[
G = \cup_p B_{\epsilon_p}(p).
\]
Then $G$ is open. Since $p \in B_{\epsilon_p}$ this also means that $E \subseteq
G \cap Y$. Now conside a point $x \in G \cap Y$. Since $x \in G$, it belongs to
at least one open ball centred at a point in $E$. Since $x$ also belongs to $Y$,
and $E$ is open relative to $Y$, $x \in E$. Thus $x \in G \cap Y \Rightarrow
x \in E$ so that $G \cap Y \subseteq E$.
\end{proof}

\end{document}
